% Source: https://artofproblemsolving.com/community/c3h1996467_volume_and_surface_area_of_a_sphere
% Stable identity: aops-topic-1996467-q01
It is enough to show that a hemisphere of radius $R$ has volume $\dfrac{2}{3}\pi R^3$.

Compare the hemisphere (flat face down on a table) with the following solid $T$: a cylinder of radius $R$ and height $R$ standing on the table, from which a cone is removed whose apex is the centre of the bottom face and whose base is the top face of the cylinder.

\textbf{Cross-sections at height $y$} ($0\le y\le R$):
\begin{itemize}
\item The hemisphere is cut in a disc of radius $\sqrt{R^2-y^2}$ (Pythagoras), with area $\pi(R^2-y^2)$.
\item The cone has height $R$ and base radius $R$, so at height $y$ its radius is $y$. The section of $T$ is a ring between radii $y$ and $R$, with area $\pi R^2-\pi y^2$.
\end{itemize}

\begin{asy}
usepackage("amsmath");
size(8.5cm);
defaultpen(fontsize(10pt));
real R=2;
filldraw(arc((0,0),R,0,180)--cycle,paleblue,black);
real y=1.1, w=sqrt(R^2-y^2);
draw((-w,y)--(w,y),red+linewidth(1.2));
label("$\sqrt{R^2-y^2}$",(-w,y),W,red);
draw((0,0)--(0,y),dashed); label("$y$",(0,y/2),W);
draw((0,0)--(w,y),dotted); label("$R$",(w/2,y/2),SE);
label("hemisphere",(0,-0.45));
pair c=(5.5,0);
filldraw(c+(-R,0)--c+(-R,R)--c+(0,0)--cycle,paleblue,black);
filldraw(c+(R,0)--c+(R,R)--c+(0,0)--cycle,paleblue,black);
draw(c+(-R,R)--c+(R,R));
draw(c+(-R,y)--c+(-y,y),red+linewidth(1.2)); draw(c+(y,y)--c+(R,y),red+linewidth(1.2));
label("$y$",c+(-y/2,y),N); label("$R$",c+(R,R*0.3),E);
label("cylinder minus cone",c+(0,-0.45));
draw((w,y)--c+(-R,y),gray+dotted);
label("\mbox{}",(min(currentpicture,user=true)+max(currentpicture,user=true))/2,black); // keeps the text after the figure black when LaTeX typesets the labels inline
\end{asy}

The areas are equal at every height, so by Cavalieri's principle the hemisphere and $T$ have the same volume:
\[
V_{\text{hemisphere}}=\pi R^2\cdot R-\dfrac{1}{3}\pi R^2\cdot R=\dfrac{2}{3}\pi R^3 .
\]
Hence the ball has volume $\boxed{\dfrac{4}{3}\pi R^3}$. $\blacksquare$

(This is Archimedes' argument. The surface area $4\pi R^2$ is then suggested by cutting the ball into many thin pyramids with apex at the centre: their total volume is $\dfrac{1}{3}\cdot(\text{total base area})\cdot R$, which must equal $\dfrac{4}{3}\pi R^3$, giving area $4\pi R^2$; making this precise needs a limit.)
